\documentclass{tietreport}

% =========================================================
% PACKAGES
% =========================================================

\usepackage{graphicx}
\usepackage{amsmath}
\usepackage{amssymb}
\usepackage{booktabs}
\usepackage{array}
\usepackage{enumitem}
\usepackage{caption}
\usepackage{float}
\usepackage{url}
\usepackage{ragged2e}

% =========================================================
% BIBLIOGRAPHY
% =========================================================

\usepackage[
    backend=biber,
    style=numeric,
    sorting=none,
    maxnames=5,
    minnames=3,
    doi=false,
    isbn=false,
    url=true,
    eprint=false
]{biblatex}

\addbibresource{references.bib}

% =========================================================
% HYPERREF
% =========================================================

\hypersetup{
    colorlinks=false,
    hidelinks,
    pdfborder={0 0 0}
}

% =========================================================
% PAGE / FLOAT SETTINGS
% =========================================================

% Prevent unnecessary blank pages caused by openright.
\let\cleardoublepage\clearpage

% Allow figures to move naturally instead of creating
% large empty spaces.
\setcounter{topnumber}{2}
\setcounter{bottomnumber}{2}
\setcounter{totalnumber}{4}

\renewcommand{\topfraction}{0.85}
\renewcommand{\bottomfraction}{0.75}
\renewcommand{\textfraction}{0.10}
\renewcommand{\floatpagefraction}{0.70}

% Paragraph spacing
\setlength{\parindent}{0pt}
\setlength{\parskip}{0.45em}

% List spacing
\setlist[itemize]{
    leftmargin=*,
    itemsep=0.25em,
    topsep=0.3em
}

\setlist[enumerate]{
    leftmargin=*,
    itemsep=0.25em,
    topsep=0.3em
}

% =========================================================
% TITLE PAGE
% =========================================================

\TitlePageHeader{UCS 503 - Software Engineering}

\ProjectTitle{MeterIQ}

\ProjectSubTitle{
    Explainable Smart Meter Analytics Platform
}
\author{%
  \begin{tabular}{c c}
    \texttt{1024160006} & Yashit Arora \\
    \texttt{1024160008} & Aarushi Gahlawat
  \end{tabular}%
}

\TitlePageSubText{%
  Group: 3P11\\
  BE Third Year, Computer Science and Engineering
}

\AdvisorName{
    Faculty Advisor: \texttt{Mr.Jeelani}
}

\TitlePageFooterText{
    \large Mid-Semester Project Evaluation \\[0.2cm]
    \bfseries Computer Science and Engineering\\
    Department\\[0.15cm]
    Thapar Institute of Engineering and Technology, Patiala
}

% =========================================================
% DOCUMENT
% =========================================================

\begin{document}

\maketitle

% =========================================================
% CONTENTS
% =========================================================

\tableofcontents

% =========================================================
% CHAPTER 1
% =========================================================

\chapter{Introduction}

% ---------------------------------------------------------
\section{Background and Context}
% ---------------------------------------------------------

MeterIQ is an explainable smart-meter analytics platform designed to support
the analysis of electricity consumption data. The system focuses on two major
analytical tasks: identifying unusual consumption behaviour and performing
short-term load forecasting.

The project uses smart-meter data from the Mathura and Bareilly regions.
Smart-meter readings contain temporal consumption information that can be
analysed to understand normal usage patterns, unusual observations, and
future demand.

The platform is designed as a decision-support system for authorised users.
Instead of automatically declaring an observation as electricity theft, the
system is intended to identify anomalous behaviour and provide risk-oriented
information that can support further human investigation.

The project combines data preparation, exploratory data analysis, feature
engineering, machine-learning analysis, forecasting, and an interactive
dashboard into a unified software system.

% ---------------------------------------------------------
\section{Problem Statement}
% ---------------------------------------------------------

Traditional examination of smart-meter data can require manual inspection of
large volumes of consumption records. This makes it difficult to quickly
identify unusual consumption patterns and understand their possible causes.

There is therefore a need for a software platform that can:

\begin{itemize}
    \item organise and process smart-meter consumption data;
    \item identify unusual consumption behaviour;
    \item provide short-term load forecasting;
    \item present analytical results in an understandable form;
    \item provide explanations and contributing factors for flagged observations;
    \item support authorised personnel during investigation; and
    \item provide a unified dashboard for accessing the analytical results.
\end{itemize}

MeterIQ addresses these requirements by combining data analytics and
visualisation in a single software-engineering system.

% ---------------------------------------------------------
\section{Project Objectives}
% ---------------------------------------------------------

The main objectives of MeterIQ are:

\begin{enumerate}
    \item Understand and prepare smart-meter consumption datasets.
    \item Perform exploratory data analysis to identify consumption patterns
    and unusual observations.
    \item Develop suitable temporal and consumption-based features.
    \item Develop an anomaly and electricity-theft-risk analysis module.
    \item Develop a short-term load forecasting module.
    \item Provide explainable analytical results rather than only binary
    decisions.
    \item Develop an interactive dashboard for authorised users.
    \item Integrate the analytical modules into a unified software system.
\end{enumerate}

% =========================================================
% CHAPTER 2
% =========================================================

\chapter{Methodology and Design}

% ---------------------------------------------------------
\section{System Architecture}
% ---------------------------------------------------------

MeterIQ follows a modular architecture in which smart-meter data is first
prepared and analysed before being passed to the analytical modules and
dashboard.

The major processing stages are:

\begin{enumerate}
    \item Smart-meter data input.
    \item Data preprocessing.
    \item Exploratory data analysis.
    \item Feature engineering.
    \item Anomaly and electricity-theft-risk analysis.
    \item Short-term load forecasting.
    \item Explainable dashboard.
\end{enumerate}

The architecture allows the individual modules to be developed and tested
independently before complete system integration.

% ---------------------------------------------------------
\subsection{High-Level Design}
% ---------------------------------------------------------

The high-level workflow of MeterIQ can be represented as:

\begin{center}
\textbf{
Smart-Meter Data
$\rightarrow$
Data Preparation
$\rightarrow$
Feature Engineering
$\rightarrow$
Analytics
$\rightarrow$
Dashboard
}
\end{center}

The analytics stage contains two major components:

\begin{itemize}
    \item anomaly and electricity-theft-risk analysis; and
    \item short-term load forecasting.
\end{itemize}

The dashboard is intended to present the results of both components through
an interactive interface.

% ---------------------------------------------------------
\subsection{System Design}
% ---------------------------------------------------------

The system is designed around four major actors or functional areas:

\begin{itemize}
    \item smart-meter data source;
    \item data and analytics processing;
    \item authorised user; and
    \item dashboard and decision-support interface.
\end{itemize}

The system design diagrams described in the following sections provide
different views of the same platform.

% =========================================================
% SOFTWARE ENGINEERING DESIGN
% =========================================================

\section{Software Engineering Design}

The software engineering design of MeterIQ defines the interaction between
users and the system, the flow of information through the platform, the
underlying entities, and the overall operational workflow.

% ---------------------------------------------------------
\subsection{Use Case Diagram}
% ---------------------------------------------------------

The Use Case Diagram represents the interaction between the authorised user
and the MeterIQ system.

The authorised user can access smart-meter information, analyse consumption,
review anomaly and electricity-theft-risk indicators, inspect explanations,
view forecasts, and use the dashboard for operational decision support.
\begin{figure}[htbp]
    \centering
    \includegraphics[
        width=\textwidth,
        height=0.62\textheight,
        keepaspectratio,
        trim=2.5cm 4.0cm 2.5cm 4.0cm,
        clip
    ]{MeterIQ_Use_Case_Diagram.pdf}
    \caption{MeterIQ Use Case Diagram}
    \label{fig:usecase}
\end{figure}

% ---------------------------------------------------------
\subsection{Data Flow Diagram}
% ---------------------------------------------------------

The Data Flow Diagrams describe how smart-meter data moves through the
MeterIQ system and how analytical results are delivered to users.

\subsubsection{DFD Level 0}

The Level 0 DFD provides a context-level representation of the MeterIQ
system.

\begin{figure}[htbp]
    \centering
    \includegraphics[
        width=0.95\textwidth,
        height=0.50\textheight,
        keepaspectratio
    ]{dfd_level0.png}
    \caption{MeterIQ Data Flow Diagram Level 0}
    \label{fig:dfd0}
\end{figure}

\subsubsection{DFD Level 1}

The Level 1 DFD decomposes the MeterIQ system into major processing stages,
including data preparation, analytics, and presentation.

\begin{figure}[htbp]
    \centering
    \includegraphics[
        width=0.95\textwidth,
        height=0.58\textheight,
        keepaspectratio
    ]{dfd_level1.png}
    \caption{MeterIQ Data Flow Diagram Level 1}
    \label{fig:dfd1}
\end{figure}

\subsubsection{DFD Level 2}

The Level 2 DFD provides a more detailed representation of the analytical
processing performed within the system.

\begin{figure}[htbp]
    \centering
    \includegraphics[
        width=0.95\textwidth,
        height=0.58\textheight,
        keepaspectratio
    ]{MeterIQ_DFD_Level2.png}
    \caption{MeterIQ Data Flow Diagram Level 2}
    \label{fig:dfd2}
\end{figure}

% ---------------------------------------------------------
\subsection{Entity Relationship Diagram}
% ---------------------------------------------------------

The Entity Relationship Diagram represents the major entities and their
relationships within the MeterIQ system.

The database design supports smart-meter information, consumption records,
analytical outputs, anomaly results, risk information, and
investigation-related information.

\begin{figure}[htbp]
    \centering
    \includegraphics[
        width=0.95\textwidth,
        height=0.58\textheight,
        keepaspectratio
    ]{ERD_MeterIQ.drawio (1).png}
    \caption{MeterIQ Entity Relationship Diagram}
    \label{fig:erd}
\end{figure}

% ---------------------------------------------------------
\subsection{Activity Diagram}
% ---------------------------------------------------------

The Activity Diagram represents the sequence of activities involved in
processing smart-meter information through the MeterIQ platform.

The general workflow starts with smart-meter data, followed by data
preparation and feature generation. The prepared information is then used
by the analytical modules, after which the results are presented to the
authorised user.

\begin{figure}[htbp]
    \centering
    \includegraphics[
        width=0.95\textwidth,
        height=0.60\textheight,
        keepaspectratio
    ]{MeterIQ_Activity_Diagram_NEW (1).drawio.png}
    \caption{MeterIQ Activity Diagram}
    \label{fig:activity}
\end{figure}

% =========================================================
% TECHNICAL STACK
% =========================================================

\section{Technical Stack}

The project uses technologies appropriate for data analytics, machine
learning, and interactive dashboard development.

\begin{table}[htbp]
\centering
\caption{MeterIQ Technical Stack}
\label{tab:techstack}
\begin{tabular}{>{\bfseries}p{0.28\textwidth} p{0.62\textwidth}}
\toprule
Area & Technology / Tool \\
\midrule
Programming & Python, JavaScript \\
Data Analysis & Pandas, NumPy \\
Machine Learning & Scikit-learn \\
Visualisation & Matplotlib, Recharts \\
Frontend & React and Vite \\
Development & VS Code, Jupyter / Google Colab \\
Version Control & Git and GitHub \\
Dashboard & React-based interactive dashboard \\
\bottomrule
\end{tabular}
\end{table}

% =========================================================
% IMPLEMENTATION DETAILS
% =========================================================

\section{Implementation Details}

% ---------------------------------------------------------
\subsection{Data Preprocessing}
% ---------------------------------------------------------

The smart-meter data is first examined and prepared before analytical
processing.

The preprocessing stage focuses on:

\begin{itemize}
    \item checking the structure of the input datasets;
    \item handling date and time information;
    \item identifying duplicate observations;
    \item checking missing or inconsistent values;
    \item maintaining appropriate temporal ordering; and
    \item preparing the data for subsequent feature engineering.
\end{itemize}

The preprocessing stage is important because incorrect temporal ordering or
data leakage can produce misleading analytical results.

% ---------------------------------------------------------
\subsection{Exploratory Data Analysis}
% ---------------------------------------------------------

Exploratory Data Analysis is used to understand the characteristics of
smart-meter consumption data before model development.

The analysis focuses on:

\begin{itemize}
    \item consumption distributions;
    \item temporal consumption patterns;
    \item hourly and daily behaviour;
    \item unusual consumption observations;
    \item variability in consumption; and
    \item relationships between temporal and consumption variables.
\end{itemize}

EDA helps establish an understanding of normal consumption behaviour before
developing analytical models.

% ---------------------------------------------------------
\subsection{Feature Engineering}
% ---------------------------------------------------------

Temporal and lag-based features are created from the prepared smart-meter
data.

The planned feature set includes:

\begin{itemize}
    \item hour;
    \item day;
    \item month;
    \item weekday;
    \item lagged consumption;
    \item previous-day consumption; and
    \item previous-week consumption.
\end{itemize}

The same-hour-last-week concept is particularly relevant to the forecasting
module because electricity demand can exhibit recurring weekly patterns.

Feature engineering is performed while maintaining chronological ordering
to reduce the possibility of data leakage.

% ---------------------------------------------------------
\subsection{Module 1: Anomaly and Electricity-Theft-Risk Analysis}
% ---------------------------------------------------------

The first analytical module focuses on identifying unusual smart-meter
consumption behaviour.

The purpose of this module is not to automatically declare a meter as
committing electricity theft. Instead, unusual observations are converted
into risk-oriented analytical information that can support investigation by
authorised personnel.

The module is expected to consider:

\begin{itemize}
    \item unusual consumption patterns;
    \item deviation from expected behaviour;
    \item temporal anomalies;
    \item relevant meter-level information;
    \item anomaly scores; and
    \item contributing factors that explain why an observation was flagged.
\end{itemize}

The final implementation and validation of the anomaly-detection component
are part of the subsequent development stage.

% ---------------------------------------------------------
\subsection{Module 2: Short-Term Load Forecasting}
% ---------------------------------------------------------

The second analytical module focuses on short-term load forecasting.

The forecasting task is designed to estimate future electricity consumption
for short horizons ranging from approximately 15 minutes to 24 hours.

A same-hour-last-week baseline is used as an important reference point.
Forecasting performance will be evaluated using suitable error metrics such
as Mean Absolute Error (MAE) and Mean Absolute Percentage Error (MAPE).

The evaluation will use a chronological holdout rather than a random split
so that future observations do not influence training information.

% ---------------------------------------------------------
\subsection{Module 3: Explainable Dashboard}
% ---------------------------------------------------------

The dashboard provides a unified interface for accessing smart-meter
analytics.

The dashboard is designed to provide:

\begin{itemize}
    \item meter search and selection;
    \item consumption history;
    \item anomaly and risk information;
    \item contributing factors and explanations;
    \item short-term forecasts;
    \item visual analytics; and
    \item operational decision-support information.
\end{itemize}

The initial frontend implementation has been developed using React, Vite,
and Recharts. At the current stage, the dashboard provides the interface
foundation, while complete backend and machine-learning integration remains
part of the next development stage.

% ---------------------------------------------------------
\subsection{Backend and Database Development}
% ---------------------------------------------------------

The backend and persistent database layer are planned for the next stage of
development.

The database will provide persistent storage for relevant meter information,
consumption records, analytical outputs, risk information, forecast outputs,
and investigation-related information.

The backend and database will subsequently be integrated with the existing
frontend and machine-learning modules.

% =========================================================
% CHAPTER 3
% =========================================================

\chapter{Results and Evaluation}

% ---------------------------------------------------------
\section{Current Project Status}
% ---------------------------------------------------------

The current mid-semester development has established the software-engineering
and data-analysis foundation of the MeterIQ platform.

The following activities have been completed or substantially developed:

\begin{enumerate}
    \item Project planning and Gantt chart preparation.
    \item Requirements and system-level design.
    \item Use Case Diagram.
    \item Data Flow Diagrams.
    \item Exploratory Data Analysis.
    \item Data preprocessing.
    \item Feature engineering.
    \item Prepared dataset validation.
    \item Initial dashboard development.
\end{enumerate}

The remaining development activities include complete machine-learning
implementation, forecasting validation, backend development, persistent
database integration, complete system integration, and final testing.

% ---------------------------------------------------------
\section{Testing Methodology}
% ---------------------------------------------------------

Testing will be performed at multiple levels to verify the correctness and
reliability of the MeterIQ system.

\subsection{Data Testing}

Data testing will verify:

\begin{itemize}
    \item input structure;
    \item missing values;
    \item duplicate records;
    \item temporal ordering;
    \item feature consistency; and
    \item processed dataset validity.
\end{itemize}

\subsection{Model Testing}

The anomaly and forecasting modules will be evaluated using appropriate
datasets and evaluation procedures.

For anomaly analysis, the evaluation will focus on whether unusual
behaviour can be identified consistently and whether the generated
explanations correspond to the underlying features.

For forecasting, chronological validation will be used to measure forecasting
errors without introducing future information into the training stage.

\subsection{Interface Testing}

The dashboard will be tested for:

\begin{itemize}
    \item navigation;
    \item meter selection;
    \item visualisation rendering;
    \item display of analytical results;
    \item explanation presentation; and
    \item usability of the main workflow.
\end{itemize}

\subsection{Integration Testing}

Integration testing will verify the complete flow:

\begin{center}
\textbf{
Data
$\rightarrow$
Processing
$\rightarrow$
Analytics
$\rightarrow$
Backend
$\rightarrow$
Dashboard
}
\end{center}

% ---------------------------------------------------------
\section{Performance Metrics}
% ---------------------------------------------------------

The final evaluation will use metrics appropriate to each analytical
component.

For forecasting, the primary metrics will include:

\begin{itemize}
    \item Mean Absolute Error (MAE);
    \item Mean Absolute Percentage Error (MAPE); and
    \item comparison against the same-hour-last-week baseline.
\end{itemize}

For anomaly analysis, evaluation will depend on the availability and quality
of appropriate validation information. The final metrics will therefore be
reported after the anomaly-detection implementation and validation have been
completed.

No final model accuracy values are reported at the current stage because the
machine-learning modules are still under development.

% ---------------------------------------------------------
\section{Analysis}
% ---------------------------------------------------------

The current implementation demonstrates the foundation required for an
integrated smart-meter analytics platform.

The project has progressed from data understanding and preparation toward
the development of analytical and presentation components.

An important design principle is explainability. Instead of presenting only
a binary decision, the system is intended to present risk-oriented
information together with contributing factors so that authorised users can
perform further investigation.

The final analysis will be completed after the anomaly-detection,
forecasting, backend, and integration stages have been implemented and
validated.

% =========================================================
% CHAPTER 4
% =========================================================

\chapter{Conclusion}

% ---------------------------------------------------------
\section{Summary of Work}
% ---------------------------------------------------------

MeterIQ is being developed as an explainable smart-meter analytics platform
for supporting distribution-grid monitoring and analysis.

The project combines smart-meter data preparation, exploratory data
analysis, feature engineering, anomaly and electricity-theft-risk analysis,
short-term load forecasting, and an interactive dashboard within a unified
software system.

At the current stage, the major software-engineering diagrams, data
exploration, preprocessing, feature engineering, prepared dataset
validation, project planning, and initial dashboard implementation have been
completed.

% ---------------------------------------------------------
\section{Key Findings}
% ---------------------------------------------------------

The current development has established the following:

\begin{itemize}

    \item \textbf{Unified Analytics:}
    Different smart-meter analytical functions can be organised within a
    single software platform.

    \item \textbf{Explainability:}
    Anomaly analysis can be presented using risk-oriented information and
    contributing factors rather than only a binary decision.

    \item \textbf{Human-in-the-Loop:}
    Electricity-theft-risk signals are intended to support further
    investigation by authorised personnel rather than automatically
    determining theft.

    \item \textbf{Temporal Analysis:}
    Smart-meter data contains temporal patterns that are useful for both
    anomaly analysis and forecasting.

    \item \textbf{Dashboard Integration:}
    An interactive dashboard provides a common interface through which
    different analytical outputs can eventually be accessed.

\end{itemize}

% ---------------------------------------------------------
\section{Future Work and Recommendations}
% ---------------------------------------------------------

The planned future work includes:

\begin{enumerate}

    \item Complete implementation and validation of the anomaly-detection
    module.

    \item Develop and evaluate short-term load forecasting models.

    \item Compare forecasting models against the same-hour-last-week
    baseline.

    \item Integrate the machine-learning modules with the dashboard.

    \item Implement backend services.

    \item Implement persistent database storage.

    \item Connect processed smart-meter datasets to the complete system.

    \item Complete end-to-end system testing.

    \item Improve explanation and investigation workflows.

    \item Evaluate the final system using appropriate performance metrics.

    \item Further improve the analytical modules based on evaluation results.

\end{enumerate}

% ---------------------------------------------------------
\section{Final Remarks}
% ---------------------------------------------------------

MeterIQ aims to provide an integrated and explainable approach to smart-meter
analytics.

The system is designed to assist authorised operators by highlighting
unusual patterns, providing forecasting information, and presenting
analytical evidence through an interactive interface.

The current mid-semester implementation establishes the foundation for the
remaining machine-learning, backend, database, integration, and testing
stages.

% =========================================================
% REFERENCES
% =========================================================

\clearpage

\printbibliography[title={References}]

% =========================================================
% END DOCUMENT
% =========================================================

\end{document}
